\begin{table}[htbp]
\centering
\footnotesize
\caption{Teste do modelo final nos quadros válidos pelo critério Q}
\label{tab:modelo-final}
\begin{tabular}{lrrrrrr}
\toprule
Modelo & F1 evento & Acurácia bal. & Macro F1 quadros & Macro F1 anotados & F1 evento com pós & Dependência \\
\midrule
Modelo do TCC (rótulos por intervalo) & 0,017 & 0,574 & 0,392 & 0,328 & 0,162 & 0,235 \\
Final da rodada 2 (janela de 2 s) & 0,046 $\pm$ 0,011 & 0,595 $\pm$ 0,018 & 0,462 $\pm$ 0,008 & 0,490 $\pm$ 0,010 & 0,219 $\pm$ 0,025 & 0,056 $\pm$ 0,015 \\
\textbf{Modelo final} (janela de 4 s) & 0,064 $\pm$ 0,005 & 0,584 $\pm$ 0,015 & 0,458 $\pm$ 0,018 & 0,510 $\pm$ 0,057 & 0,272 $\pm$ 0,014 & 0,069 $\pm$ 0,041 \\
\bottomrule
\end{tabular}
\legend{Predições fora da partição, 90.722 quadros válidos pelo critério Q comuns aos três modelos; agregado sobre os vídeos; média $\pm$ desvio entre 3 sementes (o modelo do TCC é uma execução). F1 evento: IoU 0,3, regra principal. Pós-processamento: mediana de 2 s e junção de trechos separados por menos de 1 s, escolhidos na validação. Dependência: macro F1 nos quadros anotados com as janelas reais menos com as janelas congeladas (postura média, deslocamentos zerados). Fonte: Elaborado pelo autor.}
\end{table}
